%! AP Statistics — Unit 2, Lesson 2.2 — Homework
\documentclass[10pt]{article}
\usepackage{apstats-article}
\usepackage{apstats-boxes}

%=============================================================================
\begin{document}

\pageheader{Unit 2, Lesson 2.2}{Homework}
\namedateperiod

\vspace{0.1in}

\begin{tcolorbox}[colback=sky, colframe=navy, boxrule=0.8pt, arc=2mm,
    left=3mm, right=3mm, top=1.5mm, bottom=1.5mm]
\small For each problem: (a) complete any missing table values, (b) compute conditional
relative frequencies, and (c) state whether there is evidence of an association.
Justify each conclusion with specific values.
\end{tcolorbox}

\vspace{0.14in}

% ════════════════════════════════════════════════════════════════════════════
% PROBLEM 1
% ════════════════════════════════════════════════════════════════════════════
\begin{blurbbox}[Problem 1]{navylight}

\small\textbf{Context:} 200 employees at two companies (Nexus / Apex) were surveyed
about whether they work remotely full-time (Remote: Yes/No).

\vspace{8pt}
\renewcommand{\arraystretch}{1.8}
\begin{tabularx}{\linewidth}{@{} l C{2.4cm} C{2.4cm} C{2.2cm} @{}}
\rowcolor{sky}
                     & \textbf{Remote: Yes} & \textbf{Remote: No} & \textbf{Row Total} \\
\hline
\textbf{Nexus}       & 60  & 40  & 100 \\
\textbf{Apex}        & 30  & 70  & 100 \\
\hline
\textbf{Col.\ Total} & \blank{1.6cm} & \blank{1.6cm} & \blank{1.6cm} \\
\end{tabularx}

\vspace{10pt}
\textbf{(a)} Fill in the column totals and grand total.

\vspace{8pt}
\textbf{(b)} Compute the conditional relative frequency of working remotely for each company:

\vspace{6pt}
P(Remote $|$ Nexus) $= \blank{1.2cm} / \blank{1.2cm} = \blank{2cm}$
\hfill
P(Remote $|$ Apex) $= \blank{1.2cm} / \blank{1.2cm} = \blank{2cm}$

\vspace{8pt}
\textbf{(c)} Is there evidence of an association between company and remote work status?
Justify.\\[6pt]
\writeline\\[1pt]\writeline

\end{blurbbox}

\vspace{0.14in}

% ════════════════════════════════════════════════════════════════════════════
% PROBLEM 2
% ════════════════════════════════════════════════════════════════════════════
\begin{blurbbox}[Problem 2]{greenacc}

\small\textbf{Context:} 300 high school students in a survey reported their sleep
duration category (Short: $<$7 hr / Adequate: 7--9 hr) and their self-rated academic
stress (Low / Moderate / High).

\vspace{8pt}
\renewcommand{\arraystretch}{1.8}
\begin{tabularx}{\linewidth}{@{} l C{2.0cm} C{2.0cm} C{2.0cm} C{2.0cm} @{}}
\rowcolor{sky}
                      & \textbf{Low} & \textbf{Moderate} & \textbf{High} & \textbf{Row Total} \\
\hline
\textbf{Short sleep}  & 18 & 42 & 90 & \blank{1.6cm} \\
\textbf{Adequate}     & 60 & 54 & 36 & \blank{1.6cm} \\
\hline
\textbf{Col.\ Total}  & \blank{1.4cm} & \blank{1.4cm} & \blank{1.4cm} & \blank{1.4cm} \\
\end{tabularx}

\vspace{10pt}
\textbf{(a)} Fill in all missing totals.

\vspace{8pt}
\textbf{(b)} Compute the conditional distribution of academic stress for each sleep group:

\vspace{6pt}
\renewcommand{\arraystretch}{1.8}
\begin{tabularx}{\linewidth}{@{} l C{2.6cm} C{2.6cm} C{2.6cm} C{1.6cm} @{}}
\rowcolor{sky}
                     & \textbf{Low} & \textbf{Moderate} & \textbf{High} & \textbf{Total} \\
\hline
Short sleep & \blank{2.2cm} & \blank{2.2cm} & \blank{2.2cm} & 100\% \\
Adequate    & \blank{2.2cm} & \blank{2.2cm} & \blank{2.2cm} & 100\% \\
\end{tabularx}

\vspace{10pt}
\textbf{(c)} In the space below, sketch a \textbf{segmented bar chart}.
Label both axes and include a legend.

\vspace{0.75in}

\textbf{(d)} State whether there is evidence of an association between sleep duration
and academic stress. Justify with two specific conditional percentages.\\[6pt]
\writeline\\[1pt]\writeline\\[1pt]\writeline

\end{blurbbox}

\vspace{0.14in}

% ════════════════════════════════════════════════════════════════════════════
% PROBLEM 3
% ════════════════════════════════════════════════════════════════════════════
\begin{blurbbox}[Problem 3]{redacc}

\small\textbf{Context:} A student surveyed 150 classmates on two categorical variables:
whether they prefer reading print books or e-books, and whether they prefer listening to
music while studying. Here is the \emph{joint relative frequency} table:

\vspace{8pt}
\renewcommand{\arraystretch}{1.8}
\begin{tabularx}{\linewidth}{@{} l C{2.6cm} C{2.6cm} C{2.4cm} @{}}
\rowcolor{sky}
                     & \textbf{Music: Yes} & \textbf{Music: No} & \textbf{Row Total} \\
\hline
\textbf{Print books} & 0.24 & 0.16 & \blank{1.8cm} \\
\textbf{E-books}     & 0.36 & 0.24 & \blank{1.8cm} \\
\hline
\textbf{Col.\ Total} & \blank{1.6cm} & \blank{1.6cm} & \blank{1.6cm} \\
\end{tabularx}

\vspace{10pt}
\textbf{(a)} Fill in all missing row and column totals. Verify they sum to 1.00.

\vspace{8pt}
\textbf{(b)} Compute the conditional relative frequency of listening to music for
each book-preference group:

\vspace{6pt}
P(Music $|$ Print) $= \blank{1.2cm} / \blank{1.2cm} = \blank{2cm}$
\hfill
P(Music $|$ E-books) $= \blank{1.2cm} / \blank{1.2cm} = \blank{2cm}$

\vspace{8pt}
\textbf{(c)} Is there evidence of an association? Justify.\\[6pt]
\writeline\\[1pt]\writeline

\end{blurbbox}

\vspace{0.14in}

% ── REFLECTION ───────────────────────────────────────────────────────────────
\begin{reflectionbox}
\small \textbf{Reflection:} Why do we use \emph{conditional} relative frequencies
(rather than joint relative frequencies or raw counts) to determine whether two
categorical variables are associated? Use one of the problems above as an example.

\vspace{6pt}
\writeline\\[1pt]\writeline\\[1pt]\writeline\\[1pt]\writeline
\end{reflectionbox}

\vspace{0.14in}

% ── EXTENSION ────────────────────────────────────────────────────────────────
\begin{extensionbox}
\small
\textbf{Real Data --- Pew Research}\\[4pt]
The Pew Research Center publishes tables showing survey results broken down by
demographic variables (e.g., social media use by age group).
\begin{enumerate}[leftmargin=1.5em, itemsep=8pt]
  \item Find a two-way table from a published source (or use one provided by your teacher).
        Record the source, the two variables, and the group sizes.\\[4pt]
        \writeline\\[1pt]\writeline

  \item Compute the conditional relative frequencies for at least two groups and describe
        any association you find.\\[4pt]
        \writeline\\[1pt]\writeline

  \item \textbf{Preview:} In Lesson 2.3 we will put a number on the association you described
        --- using the difference in proportions and relative risk. Write a question you have about
        how to quantify ``how strong'' the association is.\\[4pt]
        \writeline\\[1pt]\writeline
\end{enumerate}
\end{extensionbox}

\end{document}
